% Chapter 4a introduction --- scope per Project_specification.md \S4.1.
% Absorbs the killing/catastrophe distinction from the former
% 02_specification_of_rupture_state.tex (typos fixed).

\section{Introduction}

The standard models of within-host infection treat the infected cell as a
black box. In the Basic Model of Viral Replication, a cell that has been
infected releases new particles at a constant rate $p$ until it is removed
at a constant rate $d_{\Icell}$; nothing about the cell's interior enters
the dynamics, and the two constants are phenomenological --- parameters to
be fitted, not quantities to be derived. For a budding virus this is at
least a defensible caricature. For a lytic pathogen it is plainly wrong in
mechanism, whatever it may achieve by fitting: such a pathogen grows inside
the cell, and releases nothing until the cell ruptures, at which point it
releases everything at once. \textit{Yersinia pestis} in the macrophage,
\textit{Francisella tularensis} in the macrophage, the lytic phage in its
bacterium --- all follow the intracellular replication-and-lysis picture,
and all are poorly served by a constant drip.

Any replacement of that black box must begin with a complete description of
a single infected cell, and that description is the subject of this
chapter. The intracellular model was set up in Chapter~3: the
birth--death--catastrophe process, in which the pathogen count within one
cell is born at rate $\lambda$ per particle, dies at rate $\mu$ per
particle, and triggers the rupture of the cell at rate $\delta$ per
particle. Chapter~3 derived the principal single-cell quantities --- the
no-rupture probability $I(t)$, the non-fixation probability $\Ifix(t)$,
the mean and second moment $J(t)$, $K(t)$, and the cumulative mean release
$V(t)$ --- each in closed form. What it did not derive is the full
distribution of the process, or the statistics of the rupture event itself.
These are the subjects of the present chapter.

Three questions organise it.
\begin{enumerate}
\item What is the complete single-cell description? The state
probabilities at every time, the quasi-stationary distribution of the
intracellular load, and the laws of the burst time and the burst size
(\S\S\ref{sec:pgf-pdes}--\ref{sec:cond-means}).
\item How do several founders in one cell behave? The multi-founder
quantities foreshadowed in Chapter~3, and the conditional means available
to an experiment that observes a burst
(\S\ref{sec:cond-means}--\S\ref{sec:moi}).
\item What happens when the released particles found a new infection
immediately? The chained immediate-transfer model, in which one cell's
rupture is the next cell's inoculum (\S\ref{sec:chained}).
\end{enumerate}

The answers, in advance and in brief, are these. The single-cell process
is solvable far beyond its moments: the state probabilities are geometric
at every time, and their conditional limit --- the quasi-stationary
distribution --- is geometric with ratio $1/a$. The burst-size
distribution, conditioned on bursting, is \emph{exactly the same law}: a
cell's rupture size is distributed as the load it would have had, had it
kept growing forever. Burst times given burst sizes obey a harmonic-number
law; cells infected by $k$ particles release superlinearly more, and the
chained transfer of bursts from cell to cell has an exact negative-binomial
structure. Each of these results is stated in Chapter~3's notation, which
this chapter completes rather than revises.

The chapter that follows this one, Chapter~4b, takes the completed
single-cell theory and embeds it in a population: the survival and release
kernels derived here become the age-structured kernels of a renewal model
of viral dynamics, and the constant-rate black box of the Basic Model is
shown to be a projection of that model, with quantifiable distortions.
Nothing beyond the contents of the present chapter is needed for that
step.

\subsection{Two varieties of rupture: killing and catastrophe}
\label{sec:rupture-conventions}

Before we begin, it is important to state a distinction between two
varieties of the model, already touched on in Chapter~3. When a rupture
event occurs, the process goes to a state. In the literature, this state is
sometimes chosen to be $0$, and in other cases there is a separate rupture
state defined, sometimes labelled $H$. The analysis does differ between the
two options. And so for clarity, in this chapter we will refer to the
$0$-state case as ``killing''; the distinct rupture state as
``catastrophe''. The distinction matters most in
\S\S\ref{sec:pgf-pdes}--\ref{sec:state-probs}, where both conventions are
carried in parallel; from \S\ref{sec:qsd} onward the catastrophe
convention is primary.
